\documentclass{article}

\usepackage[a4paper, left=1.5cm, right=1.5cm, top=2cm, bottom=2cm]{geometry}

\usepackage{../../../../components/components}

\usepackage{hyperref}
\usepackage{fancyhdr}


% Configuration des en-têtes et pieds de page
\pagestyle{fancy}
\fancyhf{} % reset tout

\fancyhead[L]{DL2 Math-Info PR4}
\fancyhead[C]{Probabilités}
\fancyhead[R]{2025-2026}

\fancyfoot[L]{Ewen Rodrigues de Oliveira}
\fancyfoot[R]{\thepage}

\begin{document}

\docTitle{Chapitre 8 : Chaînes de Markov}

\section{Processus stochastique et propriété de Markov}
\definition{
    Soit $(\Omega, \mathcal{P})$ un espace probabilisé et soit $E$ un espace dénombrable.\\
    Une suite $(X_n)_{n \in \mathbb{N}}$ de variables aléatoires à valeurs dans $E$ est appelée un \textbf{processus stochastique} à valeurs dans $E$.
}

\example{
    On peut appliquer ça à l'évolution du capital d'un joueur dans un casino, du stock de baguettes d'une boulangerie ou encore à la météo (soleil, pluie, nuageux, etc.).
}
\definition{
    Une \textbf{chaîne de Markov} à valeurs dans un ensemble dénombrable $E$ est un processus stochastique $(X_n)_{n\in\mathbb{N}}$ à temps discret tel que pour tout $n \in \mathbb{N}$ et pour tout $(i_0, \ldots, i_{n-1}, i, j) \in E^{n+2}$ :
    \[
        \mathbb{P}(X_{n+1} = j \mid X_n = i, X_{n-1} = i_{n-1}, \ldots, X_0 = i_0) = \mathbb{P}(X_{n+1} = j \mid X_n = i)
    \]
}

\example{
    Soit $(\varepsilon_n)$ une suite de variables aléatoires indépendantes telles que $\mathbb{P}(\varepsilon_n = \pm 1) = \frac{1}{2}$.\\
    $S_{n+1} = S_n + \varepsilon_{n} = a \in \mathbb{Z}$ est une chaîne de Markov.
}

\definition{
    Une chaîne de Markov est dite \textbf{homogène} si :
    \[
        \mathbb{P}(X_{n+1} = j | X_n = i) = \mathbb{P}(X_1 = j | X_0 = i) := p_{ij} \; \forall n \in \mathbb{N}, i,j \in E
    \]
}
\vocabulary{
    On dit que $p_{ij}$ est la \textbf{probabilité de transition} entre l'état $i$ et $j$.
}
\remark{$\sum_{j \in E} p_{ij} = 1$ pour tout $i \in E$.}

\section{Matrices de transition et probabilité $n$ pas}

\definition{
    Soit $E$ un ensemble fini.\\
    On définit la \textbf{matrice de transition} $P = (p_{ij})_{i,j \in E}$.
}

\example{
    Météo : $E = \{soleil, pluie\}$. On a :
    \begin{itemize}
        \item $P(X_1 = pluie | X_0 = soleil) = p_{sp}$
        \item $P(X_1 = soleil | X_0 = pluie) = p_{ps}$
        \item $P(X_1 = soleil | X_0 = soleil) = p_{ss}$
        \item $P(X_1 = pluie | X_0 = pluie) = p_{pp}$
    \end{itemize}
    $P = \begin{pmatrix}p_{ss} & p_{sp} \\ p_{ps} & p_{pp}\end{pmatrix} = \begin{pmatrix} 0.8 & 0.2 \\ 0.4 & 0.6 \end{pmatrix}$.
}

\definition{
    Une matrice $A=(a_{ij})_{i,j \in E}$ est dite \textbf{stochastique} si :
    \[
        \sum_{j \in E} a_{ij} = 1 \; \forall i \in E
    \]
}

\definition{
    On note $\mu_u$ la loi de $X_n$, \textit{i.e.} $\mu_u(k) = \mathbb{P}(X_n = k)$ et $\mu_n = (\mathbb{P}(X_n = k))_{k \in E}$.
}

\theorem{Proposition}{}{false}{
    Soit $E$ un ensemble fini.\\
    Alors : $\mu_{u+1} = \mu_u P$
}

\noindent{\textbf{Démonstration : }\\
    $\mathbb{P}(X_{n+1}=j)=\sum_{k\in E}\mathbb{P}(X_{n+1}=j|X_n=k)\mathbb{P}(X_n=k)=\sum_{k\in E}p_{kj} \mathbb{P}(X_n=k)=\sum_{k\in E}\mu_u(k)p_{kj}$\\
    $=(\mu_u(1), \mu_u (2), ..., \mu_n (l))\cdot    \begin{pmatrix}
        p_{11} && \cdots && p_{1f} \\ \vdots && \ddots && \vdots \\ p_{\ell1} && \cdots && p_{\ell f}
    \end{pmatrix}   =(\mu_{n+1}(1), ...,\mu_{n+1}(l))$.\\
}

\definition{
    La \textbf{matrice de transition $n$ pas} est définie par $P^{(n)}_{ij} = \mathbb{P}(X_n = j | X_0 = i)$.
}

\theorem{Proposition}{}{false}{
    Soit $E$ un ensemble fini.\\
    $P^{(n)} = P^n \Leftrightarrow P^{(n)}_{ij} = P^n_{ij} \; \forall i,j \in E$
}

\noindent{\textbf{Démonstration :} Par récurrence
    \begin{itemize}
        \item $n=0$ : $P^{(0)}=P^0=Id$.
        \item On suppose que $P{(n)}=P^n$.\\
        $P_{ij}^{(n+1)}=\mathbb{P}(X_{n+1}=j|X_0=i)
            =\sum_{k\in E}\mathbb{P}(X_{n+1}=j, X_n=k|X_0=i)$\\
            $=\sum_{k\in E}\mathbb{P}(X_{n+1}=j,X_n=k, X_0=i)/\mathbb{P}(X_0=i)
            =\sum_{k\in E}\mathbb{P}(X_{n+1}=j|X_n=k, X_0=i)\mathbb{P}(X_n=k, X_0=i)$\\
            $=\sum_{k\in E}\mathbb{P}(X_{n+1}=j|X_n=k)\mathbb{P}(X_0=i)$ 
            avec $\mathbb{P}(X_0=i)=\mathbb{P}(X_n=k|X_0=i)$\\
            $=\sum_k P_{kj} P_{ij}^{(n)}$ par chaîne de Markov.\\
            $=\sum_k P_{ik}^n P_{kj}=P_{ij}^{n+1}$ par hypothèse de récurrence.
    \end{itemize}
}

\theorem{Corollaire}{}{false}{
    Soit $\mu_n$ le vecteur ligne de probabilité à l'instant $n$, tel que $(\mu_n)_j = \mathbb{P}(X_n = j)$. 
    Alors pour tout $n \in \mathbb{N}$ :
    \[ \mu_n = \mu_0 P^n \]
}

\noindent{\textbf{Démonstration : } 
    $\mu_u(k)=\sum_{i\in E} \mathbb{P}(X_n=k|X_0=i) \mathbb{P}(X_0=i)=\sum_{i\in E} \mu_0(i) P_{ik}^n$.\\
}

\example{
    Considérons $P = \begin{pmatrix} 0.8 & 0.2 \\ 0.4 & 0.6 \end{pmatrix}$ et $\mu_0 = (1,0)$.\\
    Alors : $\mu_1 = \mu_0 P = (0.8, 0.2)$, $\mu_2 = \mu_1 P = (0.68, 0.32)$
}

\definition{
    Soit une chaîne de Markov de distribution initiale $\delta_i$ (\textit{i.e.} $\delta_i(j) = \delta_{ij}$).\\
    La \textbf{probabilité de transition à $n$ pas} de $i$ vers $j$ est donnée par $\mathbb{P}_i(X_n = j)$, où $\mathbb{P}_i$ est la loi de la chaîne de Markov.
}

\remark{On définit une matrice $P_{ij}^{(n)} = \mathbb{P}_i(X_n = j)$, qui est la matrice de transition à $n$ pas.}

\definition{
    Soit $(X_n)_{n\in\mathbb{N}}$ une chaîne de Markov à valeurs dans $E$ de matrice de transition $P$.\\
    On dit qu'une loi $\pi$ sur $E$ est une \textbf{loi stationnaire} de la chaîne de Markov si $\pi P = \pi$.
}
\vocabulary{On appelle aussi \textbf{distribution stationnaire} une loi stationnaire.}

\theorem{Théorème}{}{true}{
    Soit $E$ un ensemble fini. Pour toute chaîne de Markov à valeurs dans $E$, il existe une loi stationnaire.
}

\example{
    Soit $P = \begin{pmatrix} 0.8 & 0.2 \\ 0.4 & 0.6 \end{pmatrix}$.\\
    On cherche $\pi$ tel que $\pi P = \pi$.\\
    $\pi_1 = 0.8\pi_1 + 0.4\pi_2$ et $\pi_2 = 0.2\pi_1 + 0.6\pi_2$.\\
    En résolvant ce système, on trouve $\pi = (\frac{2}{3}, \frac{1}{3})$.
}
\example{
    On pose $E = \{0,1,2\}$ et $P = \begin{pmatrix} 0 & 1 & 0 \\ \frac{1}{2} & 0 & \frac{1}{2} \\ 0 & 0 & 1 \end{pmatrix}$.\\
    On a : $\pi = (0,0,1)$ tel que $\pi P = \pi$.
}

\section{Classes de communication}

\definition{
    Soit $E$ un espace d'état et $(X_n)_{n\in\mathbb{N}}$ une chaîne de Markov à valeurs dans $E$ dont on note $P$ la matrice de transition. Soient $i,j \in E$.\\\\
    On dit que $i$ \textbf{communique vers} $j$ (noté $i \rightarrow j$) si $\exists n \in \mathbb{N}$ tel que $p_{ij}^{(n)} > 0$.
}
\vocabulary{On dit que $i$ et $j$ \textbf{communiquent} si $i \rightarrow j$ et $j \rightarrow i$. On note alors $i \leftrightarrow j$.}

\theorem{Proposition}{}{false}{
    La relation de communication $\leftrightarrow$ est une relation d'équivalence.
}
\noindent{\textbf{Démonstration :} On montre trivialement que $\leftrightarrow$ est réflexive, symétrique et transitive.}

\vocabulary{On appelle \textbf{classes de communication} les classes d'équivalence de la relation $\leftrightarrow$.}

\definition{
    On dit d'une chaîne de Markov qu'elle est \textbf{irréductible} si elle n'a qu'une seule classe de communication, c'est-à-dire que tous les états communiquent entre eux.
}

\definition{
    Soit $i \in E$. On dit que $i$ est un \textbf{état absorbant} si $p_{ii} = 1$.
}

\definition{
    Soit $(X_n)_{n\in\mathbb{N}}$ une chaîne de Markov à valeurs dans $E$. Soit $i \in E$.\\
    On définit la \textbf{période} de $i$ par $d(i) = pgcd\{n \in \mathbb{N}^* : P_{ii}^{(n)} > 0\}$.
}
\vocabulary{On dit d'un état $i$ qu'il est \textbf{apériodique} si $d(i) = 1$.}

\remark{Si une chaîne de Markov est irréductible, alors tous les états ont la même période (qu'on appelle la période de la chaîne).}

\example{
    Posons $E = \{0,1,2\}$ et $P = \begin{pmatrix} 0 & 1 & 0 \\ \frac{1}{2} & 0 & \frac{1}{2} \\ 0 & 0 & 1 \end{pmatrix}$.
    
    L'analyse de la structure de la chaîne montre que :
    \begin{itemize}
        \item $0$ et $1$ communiquent entre eux ($0 \leftrightarrow 1$) car $P_{01} > 0$ et $P_{10} > 0$.
        \item $0$ et $1$ peuvent aller vers $2$ ($0 \to 1 \to 2$), mais le retour est impossible car $P_{20} = P_{21} = 0$.
        \item L'état $2$ est un \textbf{état absorbant} (une classe fermée à lui seul), car $P_{22} = 1$.
    \end{itemize}
    On distingue donc deux classes de communication : $\{0, 1\}$ et $\{2\}$.
    Les périodes sont : $d(0) = 2$ (car on peut revenir à $0$ en 2 étapes via $1$), $d(1) = 2$ (car on peut revenir à $1$ en 2 étapes via $0$), et $d(2) = 1$ (apériodique, car c'est un état absorbant).
}

\theorem{Théorème}{}{true}{
    Soit $(X_n)_{n\in\mathbb{N}}$ une chaîne de Markov à valeurs dans $E$. Notons $P$ sa matrice de transition.\\
    Supposons que la chaîne soit \textbf{irréductible} et \textbf{apériodique}. Alors :
    \begin{enumerate}
        \item[i)] il existe une unique loi stationnaire,
        \item[ii)] pour toute loi initiale $\mu_0$, on a $\lim_{n \to \infty} \mu_0 P^n = \pi$ (convergence vers la loi stationnaire).
    \end{enumerate}
}

\newpage
\tableofcontents

\end{document}